\chapter{Classical Information, Reversible Dynamics, and Matrices}

\lectureoverview{Quantum mechanics asks us to reason outside the range in which ordinary intuition was trained. We therefore begin one step earlier, with classical information: bits, configurations, and reversible update laws. The same discussion leads naturally to vectors and matrices, which will provide the language for quantum bits.}

\section{Intuition Built for the Wrong World}

Our physical intuition is old. Long before anyone wrote an equation, an animal chasing prey had to estimate relative speed and direction. If a lion is gaining on an antelope, the lion is responding to a changing separation and therefore, in an implicit way, to relative velocity. Nervous systems that handled such judgments well were favored; those that handled them badly did not leave many descendants. Much of what we call common sense about motion was trained in this ordinary world.

That inheritance is useful, but narrow. Nothing in ordinary life moves close to the speed of light, so evolution supplied no instinct for relativistic spacetime. A physicist first learns the abstract geometry and only later acquires a new intuition for it. Quantum mechanics is harder still. Our familiar picture of a particle comes from rocks and arrows: at each instant the object has a position, the succession of positions forms a trajectory, and changes of position define a velocity. An electron does not oblige us by fitting that picture. Position uncertainty and the absence of a classical trajectory have no analogue in the experience that shaped our senses.

There is no paradox in this failure of intuition. There would be a paradox only if instincts adapted to stones, animals, and modest speeds somehow knew in advance about electrons. The right response is to build a different intuition from a small set of precise rules. Here the emphasis will not be the usual route through wave mechanics and the Schr\"odinger equation. We will concentrate on the logical structure of quantum mechanics and on quantum information. Before asking what a quantum bit is, we must understand exactly what an ordinary bit does.

\section{Bits and the Amount of Information}

Physics is inseparable from information. Sometimes that information is expressed by a real number, but the simplest possible piece of information is an answer to a question with only two alternatives: yes or no, heads or tails, up or down.

\begin{definition}
A \emph{classical bit} is a physical question with exactly two possible answers.
\end{definition}

We may label those answers \(0\) and \(1\). The labels carry no intrinsic meaning; any two distinct symbols would do. Since Dirac notation will soon become indispensable, we introduce it even while it is redundant:
\begin{equation}
  \ket{0},
  \qquad
  \ket{1}.
\end{equation}
The vertical bar and angle bracket form a \emph{ket}, the right half of a bra--ket. Nothing quantum has entered merely because we have changed notation.

For several distinguishable bits, write the answers in a fixed order. Six distinguishable coins, for example, may have the configuration
\begin{equation}
  \ket{0\,1\,1\,0\,1\,0}.
\end{equation}
Commas are optional once the ordering is understood. One bit has two possible states, two bits have four, and every added bit doubles the count. Thus an \(n\)-bit system has
\begin{equation}
  N_s=2^n
  \label{eq:l01-state-count}
\end{equation}
states. Conversely,
\begin{equation}
  n=\log_2 N_s.
  \label{eq:l01-bits-from-states}
\end{equation}

Equation~\eqref{eq:l01-bits-from-states} suggests a definition that is useful even when the number of states is not a power of two:
\begin{equation}
  I=\log_2 N_s.
  \label{eq:l01-information}
\end{equation}
A six-sided die therefore carries \(\log_2 6\) bits of information. The answer need not be an integer. Systems assembled literally from binary switches will be convenient examples, but the logarithm is the more general measure.

\section{Encoding an Apparently Continuous World}

Bits are not confined to coins. Any measured quantity known to finite accuracy can be written as a finite binary string. Suppose, for example, that the temperature of a room is required only to a chosen number of significant figures. Express the measured value in base two and retain the digits needed for that accuracy. Greater precision requires a longer string, but no change of principle.

A spatially varying field can be treated similarly. Divide the room into a fine ordered lattice of cells, record the field value in each cell to finite precision, and concatenate the resulting binary strings. Making the cells smaller improves spatial resolution; keeping more binary digits improves the field-value resolution. Temperature fields, electric fields, and magnetic fields can all be encoded this way to any prescribed finite accuracy.

Particle configurations give an even more literal binary picture. Subdivide space into cells small enough that each contains at most one particle. Assign
\begin{equation}
  n_i=
  \begin{cases}
    1,&\text{cell \(i\) is occupied},\\
    0,&\text{cell \(i\) is empty}.
  \end{cases}
  \label{eq:l01-occupancy-bit}
\end{equation}
The complete configuration is the ordered string
\begin{equation}
  \ket{n_1n_2n_3\cdots}.
\end{equation}
It is a very long answer to a sequence of yes-or-no questions: is there a particle in the first cell, in the second, and so on? This is the physical version of the game \emph{Twenty Questions}, except that the number of questions may be enormous.

The construction explains why digital computers can simulate physics. A computer stores finite strings, while finite resolution turns measured quantities and fields into precisely such strings. The approximation can be improved systematically by refining the lattice, the time step, or the number of retained digits.

\begin{classroomqa}
\lecturetimestamp{00:29:43}
\audiencequestion{Is the accuracy of the temperature representation arbitrary?}
\lecturerresponse{It is chosen according to the approximation we want. A finite string gives finite accuracy; retaining more binary places gives a finer approximation. There is no claim here that a finite string specifies an arbitrary classical real number exactly.}
\end{classroomqa}

\begin{classroomqa}
\lecturetimestamp{00:30:11}
\audiencequestion{What happens for a general real number, especially an irrational one?}
\lecturerresponse{An exact real number may require infinitely many bits. Some rational binary expansions repeat, and a generic irrational expansion never terminates. A real computer always keeps finitely many digits. The point is that any specified finite accuracy can be encoded, not that every classical real number has a finite exact description.}
\end{classroomqa}

\section{A State and a Law of Motion}

So far we have described a system at one instant. Motion requires a second ingredient.

\begin{definition}
A \emph{configuration} is the complete state of a system at a specified instant. A discrete-time law of motion is a rule that maps each configuration to the configuration one time step later.
\end{definition}

The distinction is elementary but essential. A snapshot of an electron distribution is a configuration; it is not yet an orbit. If space is replaced by a lattice and time by a sequence \(t_0,t_1,t_2,\ldots\), then motion is the succession
\begin{equation}
  \ket{s(t_0)}
  \longmapsto
  \ket{s(t_1)}
  \longmapsto
  \ket{s(t_2)}
  \longmapsto\cdots.
  \label{eq:l01-discrete-motion}
\end{equation}
A physical model therefore needs both a set of possible states and an update rule.

\begin{classroomqa}
\lecturetimestamp{00:31:12}
\audiencequestion{How does this cell description account for an electron moving in an orbit?}
\lecturerresponse{First specify only the configuration at an instant: which cells are occupied, with the nucleus fixed if that is part of the model. Then discretize time. The orbital motion is the rule that updates one occupancy string to the next. The rule may be complicated because of the ordering of the cells, but it is still an update of information.}
\end{classroomqa}

It helps to draw an abstract \emph{state space}. Each point in this space denotes an entire configuration, not a location in ordinary space. A \(10\times10\) occupancy lattice contains one hundred bits and hence \(2^{100}\) possible configurations.\editorialnote{The spoken lecture briefly says \(10^{100}\) while developing the example; the bit count fixes the intended value as \(2^{100}\).} Each one of those configurations is a single point in the state space.

For one bit there are only two points. The identity law leaves both fixed:
\begin{equation}
  0\mapsto0,
  \qquad
  1\mapsto1.
  \label{eq:l01-identity-update}
\end{equation}
The other reversible possibility swaps them:
\begin{equation}
  0\mapsto1,
  \qquad
  1\mapsto0.
  \label{eq:l01-flip-update}
\end{equation}
This is a flip-flop. A freely evolving physical example is less immediate than a manually operated light switch, because intervention would enlarge the system. Day and night provide an approximate two-state cycle; a pendulum also alternates between extremes, although its full state includes all the intermediate positions and velocities.

\begin{classroomqa}
\lecturetimestamp{00:40:38}
\audiencequestion{Would an ordinary light switch be an example of the flip-flop law?}
\lecturerresponse{Not if a person is toggling it, because then the switch is not evolving by itself; the person belongs to the enlarged dynamical system. What we seek is a closed system whose own law alternates its state. Day and night suggest the idea, while a pendulum reminds us that a real example usually has more than two states.}
\end{classroomqa}

A two-bit system has four configurations. One update law may carry all four around a single cycle. Another may consist of two separate two-cycles. These are loops in the logical space of configurations, not necessarily loops in physical space. The arrows specify the law: wherever the system starts, following the arrows tells us its subsequent history.

\section{Reversibility and the Preservation of Distinctions}

Not every arrow diagram has the same physical status. Consider a rule in which two different states merge into the same successor:
\begin{equation}
  a\mapsto c,
  \qquad
  b\mapsto c.
  \label{eq:l01-merging-update}
\end{equation}
The next state is defined whether we start at \(a\) or \(b\), but after reaching \(c\) we cannot determine which past occurred. The distinction between \(a\) and \(b\) has disappeared.

\begin{figure}[htbp]
\centering
\includegraphics[width=0.82\linewidth]{lecture_01_figure_02.png}
\caption{Several update diagrams remain on the board during the discussion of reversibility. Closed cycles coexist with a merging pattern, and the quantum term \emph{unitarity} is written beside them.}
\lectureframe{00:53:21}
\end{figure}

\begin{figure}[htbp]
\centering
\begin{tikzpicture}[>=Latex,thick,node distance=1.15cm]
  \node[circle,fill,inner sep=1.8pt,label=above:\(1\)] (r1) {};
  \node[circle,fill,inner sep=1.8pt,right=of r1,label=above:\(2\)] (r2) {};
  \node[circle,fill,inner sep=1.8pt,below=of r2,label=below:\(3\)] (r3) {};
  \node[circle,fill,inner sep=1.8pt,left=of r3,label=below:\(4\)] (r4) {};
  \draw[->] (r1) -- (r2);
  \draw[->] (r2) -- (r3);
  \draw[->] (r3) -- (r4);
  \draw[->] (r4) -- (r1);
  \node[below=0.5cm of r4,anchor=north west] {\small reversible cycle};

  \node[circle,fill,inner sep=1.8pt,right=4.2cm of r1,label=above:\(a\)] (a) {};
  \node[circle,fill,inner sep=1.8pt,below=1.25cm of a,label=below:\(b\)] (b) {};
  \node[circle,fill,inner sep=1.8pt,right=1.6cm of a,yshift=-0.625cm,label=above:\(c\)] (c) {};
  \node[circle,fill,inner sep=1.8pt,right=of c,label=above:\(d\)] (d) {};
  \draw[->] (a) -- (c);
  \draw[->] (b) -- (c);
  \draw[->] (c) -- (d);
  \node[below=0.55cm of b,anchor=north west] {\small information-losing merge};
\end{tikzpicture}
\caption{A reversible rule assigns one predecessor and one successor to every state. A merging rule has a unique future but an ambiguous past.}
\end{figure}

Fundamental classical dynamics does not permit this loss when every degree of freedom is retained. The required property has two halves:
\begin{equation}
  \boxed{\text{one unique future and one unique past for every state}.}
  \label{eq:l01-reversibility}
\end{equation}
For a finite state space, the update must therefore be a permutation. Every state has exactly one outgoing arrow and exactly one incoming arrow, so the diagram decomposes into disjoint cycles.

\begin{classroomqa}
\lecturetimestamp{00:46:32}
\audiencequestion{What is different about the merging diagram?}
\lecturerresponse{It does not consist entirely of loops. More fundamentally, finding the system at the merged state does not tell us which predecessor it had. We can construct the future, but we cannot reconstruct the past. One bit of historical distinction has been lost.}
\end{classroomqa}

This notion is \emph{time reversibility}, not necessarily exact time-reversal symmetry. A clockwise cycle has a definite forward orientation. The original law does not send the system counterclockwise, but an inverse law exists that does. Exact time-reversal symmetry asks an additional question: after applying the appropriate reversal operation, are the same dynamical laws obtained? Reversibility alone asks only whether the history can be uniquely run backward.

\begin{classroomqa}
\lecturetimestamp{00:49:58}
\audiencequestion{Is the uniqueness of the past just time-reversal symmetry?}
\lecturerresponse{Not exactly. The arrows of a cycle may point only one way, so the law has an orientation. Reversibility means that a well-defined inverse law can be constructed by reversing every arrow. It does not by itself say that the inverse law is identical to the original law.}
\end{classroomqa}

Nor should we repair an irreversible diagram by placing two outgoing arrows at the merged state. That would make the backward motion ambiguous in a different way: the system would need a probabilistic instruction, perhaps choosing one branch \(30\%\) of the time and the other \(70\%\). We are still constructing deterministic classical dynamics, for which one present state determines one future and one past.

Every real system is ultimately quantum mechanical. Why large objects appear classical, and why their quantum alternatives become difficult to observe, is a later question. The corresponding quantum principle will preserve the same logical core. Half jokingly, one might call uniqueness of past and future \emph{unique-tarity}. Its actual name is \emph{unitarity}: quantum evolution can be run uniquely forward or backward.

\section{Apparent Information Loss and Coarse-Graining}

Macroscopic experience seems to contradict reversibility. Put all the molecules of a gas in the left corner of a sealed room and release them. Soon the gas fills the room. Start instead from the right corner and the final macroscopic picture looks the same. Have the two distinct initial states merged?

Only after we have thrown away detail. A microscopic state specifies every molecule. If every position and velocity could be followed with sufficient precision, reversing those motions would distinguish the two preparations. In practice the calculation is prohibitively difficult and the required precision grows rapidly, but practical inability is not fundamental erasure. The macroscopic statement \emph{air fills the room} identifies an enormous number of distinct microscopic configurations as one coarse-grained state.

This is where the second law of thermodynamics enters. We lose the ability to distinguish detailed states because we do not keep track of all microscopic variables. Information appears to disappear because our description is coarse, not because the fine-grained deterministic law merges distinct points.

\begin{classroomqa}
\lecturetimestamp{00:56:11}
\audiencequestion{Would tracking all the molecules in a classical room require infinitely many bits?}
\lecturerresponse{For an exact classical description, yes. Positions and velocities are real numbers, and following a chaotic many-particle system for longer times demands ever greater precision. Quantum mechanics changes this conclusion: for a sealed region with bounded energy, the physically distinguishable description can be discrete rather than an arbitrary list of exact classical real numbers.}
\end{classroomqa}

The issue is not merely academic. Whether a black hole can destroy information was disputed for decades; Stephen Hawking originally argued that it could. That controversy asks whether the apparent merging is only another coarse-grained description or a genuine violation of quantum unitarity. We will not settle it before quantum mechanics has supplied the precise language, but it reveals how much rests on the simple arrow diagrams.

\section{Rows, Columns, and Inner Products}

Before replacing classical bits by quantum bits, we need a small amount of linear algebra. For present purposes a vector may be treated simply as a finite ordered list of real numbers. A four-component vector can be written horizontally,
\begin{equation}
  A^{\mathsf T}=(A_1,A_2,A_3,A_4),
  \label{eq:l01-row-vector}
\end{equation}
or vertically,
\begin{equation}
  A=
  \begin{pmatrix}
    A_1\\ A_2\\ A_3\\ A_4
  \end{pmatrix}.
  \label{eq:l01-column-vector}
\end{equation}
The word \emph{four-dimensional} means only that there are four components. There is no need to visualize an arrow in ordinary space. The horizontal object is a row vector and the vertical object a column vector.

\begin{figure}[htbp]
\centering
\includegraphics[width=0.92\linewidth]{lecture_01_figure_03.png}
\caption{The same components are written as a row and as a column. At this stage they are real numbers, so no complex conjugation is involved.}
\lectureframe{01:03:05}
\end{figure}

\begin{classroomqa}
\lecturetimestamp{01:03:14}
\audiencequestion{Isn't the row vector the dagger of the column vector?}
\lecturerresponse{That will be the right language once the components are complex. Then transposition must be accompanied by complex conjugation. For the moment the entries are real, so we need only distinguish the row arrangement from the column arrangement.}
\end{classroomqa}

Let
\begin{equation}
  B=(B_1,B_2,B_3,B_4)
\end{equation}
be a row with the same number of entries as the column \(A\). Their product is defined by matching positions and summing:
\begin{align}
  BA
  &=
  (B_1,B_2,B_3,B_4)
  \begin{pmatrix}
    A_1\\ A_2\\ A_3\\ A_4
  \end{pmatrix}\\
  &=B_1A_1+B_2A_2+B_3A_3+B_4A_4.
  \label{eq:l01-inner-product}
\end{align}
The result is a number, not another vector or a matrix. In an abstract vector space this operation is called an inner product.

\begin{classroomqa}
\lecturetimestamp{01:06:44}
\audiencequestion{Is this the dot product?}
\lecturerresponse{Yes. For ordinary vectors in three-dimensional space it is the familiar dot product. The name \emph{inner product} emphasizes that the same row-by-column operation applies in vector spaces of any finite dimension.}
\end{classroomqa}

\section{Matrices as Machines for Updating States}

A matrix is a square array of numbers that acts on a vector to produce another vector. Think of it as a machine: feed in the old column, apply a fixed rule, and obtain a new column. In four dimensions,
\begin{equation}
  M=
  \begin{pmatrix}
    M_{11}&M_{12}&M_{13}&M_{14}\\
    M_{21}&M_{22}&M_{23}&M_{24}\\
    M_{31}&M_{32}&M_{33}&M_{34}\\
    M_{41}&M_{42}&M_{43}&M_{44}
  \end{pmatrix}.
  \label{eq:l01-general-matrix}
\end{equation}
The first index labels the row and the second the column. One may regard \(M\) as a collection of columns or as a collection of rows; both viewpoints will be useful.

To calculate \(A'=MA\), take the inner product of each row of \(M\) with the input column:
\begin{equation}
  A'_i=(MA)_i=\sum_{j=1}^{4}M_{ij}A_j.
  \label{eq:l01-matrix-vector}
\end{equation}
For example,
\begin{align}
  A'_1&=M_{11}A_1+M_{12}A_2+M_{13}A_3+M_{14}A_4,\\
  A'_2&=M_{21}A_1+M_{22}A_2+M_{23}A_3+M_{24}A_4.
\end{align}
The third and fourth entries follow by the same pattern.

\begin{figure}[htbp]
\centering
\includegraphics[width=0.94\linewidth]{lecture_01_figure_04.png}
\caption{A matrix acts on a column vector row by row. The first two output components are expanded explicitly on the board.}
\lectureframe{01:12:45}
\end{figure}

\subsection{Encoding a Classical Configuration}

Suppose a system has four possible states. Represent them by the one-hot columns
\begin{equation}
  e_1=
  \begin{pmatrix}1\\0\\0\\0\end{pmatrix},
  \quad
  e_2=
  \begin{pmatrix}0\\1\\0\\0\end{pmatrix},
  \quad
  e_3=
  \begin{pmatrix}0\\0\\1\\0\end{pmatrix},
  \quad
  e_4=
  \begin{pmatrix}0\\0\\0\\1\end{pmatrix}.
  \label{eq:l01-one-hot}
\end{equation}
The location of the single \(1\) tells us which classical configuration is present.

Consider the cyclic update
\begin{equation}
  1\longrightarrow2\longrightarrow3\longrightarrow4\longrightarrow1.
  \label{eq:l01-four-cycle}
\end{equation}
It is represented by
\begin{equation}
  P=
  \begin{pmatrix}
    0&0&0&1\\
    1&0&0&0\\
    0&1&0&0\\
    0&0&1&0
  \end{pmatrix}.
  \label{eq:l01-permutation-matrix}
\end{equation}
Indeed,
\begin{equation}
  Pe_1=e_2,\qquad
  Pe_2=e_3,\qquad
  Pe_3=e_4,\qquad
  Pe_4=e_1.
\end{equation}
The board construction is corrected while it is being written, first with five states and then with four. Equation~\eqref{eq:l01-permutation-matrix} records the settled cyclic rule rather than the abandoned intermediate array.

The same matrix acts on an arbitrary column:
\begin{equation}
  P
  \begin{pmatrix}A\\B\\C\\D\end{pmatrix}
  =
  \begin{pmatrix}D\\A\\B\\C\end{pmatrix}.
  \label{eq:l01-cyclic-shift}
\end{equation}
For a classical state only one component is \(1\), but allowing general entries exposes the underlying linear operation. Classical deterministic updates use special zero-one matrices with one \(1\) in every row and every column. Quantum mechanics will retain matrix evolution while greatly enlarging the allowed entries.

\subsection{Composition Is Matrix Multiplication}

If one time step is
\begin{equation}
  V'=MV,
\end{equation}
then a second application gives
\begin{equation}
  V''=MV'=M(MV)=M^2V.
  \label{eq:l01-two-updates}
\end{equation}
Multiplying matrices is therefore not an arbitrary algebraic pastime. It composes physical updates. Five time steps are represented by \(M^5\).

For two \(2\times2\) matrices
\begin{equation}
  M=
  \begin{pmatrix}
    M_{11}&M_{12}\\
    M_{21}&M_{22}
  \end{pmatrix},
  \qquad
  N=
  \begin{pmatrix}
    N_{11}&N_{12}\\
    N_{21}&N_{22}
  \end{pmatrix},
\end{equation}
the product \(C=MN\) has entries
\begin{equation}
  C_{ij}=\sum_k M_{ik}N_{kj}.
  \label{eq:l01-matrix-product}
\end{equation}
In particular,
\begin{align}
  C_{11}&=M_{11}N_{11}+M_{12}N_{21},\\
  C_{12}&=M_{11}N_{12}+M_{12}N_{22},\\
  C_{21}&=M_{21}N_{11}+M_{22}N_{21},\\
  C_{22}&=M_{21}N_{12}+M_{22}N_{22}.
\end{align}
Each entry is a row of the left matrix multiplied by a column of the right matrix. Equivalently, act with \(M\) on each column of \(N\) and place the resulting columns side by side.

\begin{figure}[htbp]
\centering
\includegraphics[width=0.86\linewidth]{lecture_01_figure_05.png}
\caption{Matrix multiplication organized as repeated matrix--column multiplication. The first row and the columns of the second matrix are marked to show where each entry of the product comes from.}
\lectureframe{01:26:30}
\end{figure}

A row vector can also act from the left. If \(B=(B_1,\ldots,B_4)\), then \(BM\) is another row vector with
\begin{equation}
  (BM)_j=\sum_i B_iM_{ij}.
  \label{eq:l01-row-times-matrix}
\end{equation}
Again the rule is row times column. The shapes keep track of the result:
\begin{equation}
\begin{array}{c@{\qquad}c}
\text{row}\times\text{column}&=\text{number},\\
\text{matrix}\times\text{column}&=\text{column},\\
\text{row}\times\text{matrix}&=\text{row},\\
\text{matrix}\times\text{matrix}&=\text{matrix}.
\end{array}
\label{eq:l01-shape-rules}
\end{equation}

These operations are worth practicing with small numerical \(2\times2\) and \(3\times3\) examples. Together with a little calculus, row--column multiplication is the main new mathematical habit needed for the quantum mechanics ahead.

\begin{classroomqa}
\lecturetimestamp{01:31:50}
\audiencequestion{What restriction does reversibility place on the update matrix?}
\lecturerresponse{It must have an inverse. The inverse matrix implements the reverse update and takes every state back to its unique predecessor. Not every matrix is invertible, so the algebra reproduces the distinction made earlier with arrow diagrams. For the cyclic example, \(P^{-1}=P^3\) and \(P^4=\mathbb{1}\).}
\end{classroomqa}

\section{Why Prepare the Algebra Before the Qubit?}

The first hour gave a qualitative account of deterministic classical physics in terms of bits and reversible arrows. The remaining mathematics translates that account into a form that can survive the transition to quantum mechanics. A configuration becomes a column; a one-step law becomes a matrix; repeated time evolution becomes a product of matrices; reversibility becomes the existence of an inverse.

We have not yet used a quantum rule. That delay is deliberate. Next we replace the classical alternatives \(\ket{0}\) and \(\ket{1}\) by a qubit and ask what new states are possible. The row vectors, column vectors, inner products, and matrices introduced here will then reveal why quantum information differs so sharply from a classical string of bits.

\begin{classroomqa}
\lecturetimestamp{01:34:00}
\audiencequestion{Why introduce all this matrix algebra before saying what it is for?}
\lecturerresponse{The qualitative problem came first: how should deterministic physics be represented by states and update laws? The algebra is the preparation for the next step. Once qubits are introduced, these manipulations will no longer look detached; they will be the natural language of quantum states and their evolution.}
\end{classroomqa}
